\documentclass[11pt]{article}
\usepackage[margin=0.8in]{geometry}
\usepackage{booktabs,hyperref}
\hypersetup{colorlinks=true,urlcolor=blue}
\title{Impala performance experiment: Step 2}
\author{Compiled copper material-history driver}
\date{September 29, 2026}
\begin{document}\maketitle
\section*{Outcome and scope}
The new optional Numba backend compiles the outer material-history driver for
the copper PTW model combination. The post-cache profile attributed 0.889 of
1.023 seconds (87\%) to this driver in 200 batch-50 evaluations. Individual
physics functions already used Numba; the change removes Python coordination
at each of 100 integration steps and checks invariant PTW parameter constraints
once per history call.

The kernel reuses the existing numerical functions, equations, precision and
integration resolution. It uses no fast-math, parallel execution or GPU. The
Python driver remains the default. Supported submodels are PTW stress,
quadratic specific heat, linear density, linear melt temperature and BGP/PW
shear modulus. Other combinations and overridden classes/methods are rejected
explicitly when requesting the compiled backend. All final material-state
fields are restored after the compiled call.
\section*{Accuracy checks}
All 83 tests passed in 114.72 seconds, including 17 new compiled-driver tests.
These cover pooled and hierarchical batches, changing parameters and constants,
modified conditions and integration resolution, supplied nonuniform histories,
continued nonzero states, small-p branches and strain-rate thresholds.
Predictions and all final state fields match the Python control exactly in
these comparisons. Unsupported inputs and invalid parameters are also tested.

A separate benchmark launched isolated processes importing either an untouched
original checkout, the cached Python version, or the compiled version. Each
matched all 256 original-fixture predictions and weighted log likelihoods
exactly, with maximum absolute error zero. For each variant, seeds 42 and 123
were run in pooled and hierarchical modes: 3,000 stored states, eight
temperatures, six copper experiments and 327 observations. All recorded
parameter/noise arrays and acceptance counters matched the original exactly;
hierarchical means, covariances and swap counters also matched. This independent
reference shares no modified physics modules with the candidate.

These are short-chain regression checks, not a new convergence study. No
long-run ESS-per-second experiment was repeated. The previously identified
hierarchical mixing issue remains unresolved.
\section*{Measured performance}
Each process recorded 60 warmed evaluations per batch size. The table reports
the mean of the two process medians in milliseconds. Variant order was reversed
for the second seed. Speedup here is relative to the cached Python version,
so it isolates step 2 from the earlier caching change.
\begin{center}\begin{tabular}{rrrrr}\toprule
Batch & Original (ms) & Cached (ms) & Compiled (ms) & Speedup\\\midrule
1 & 1.582 & 1.514 & 0.418 & 3.63\\
4 & 1.852 & 1.748 & 0.613 & 2.85\\
50 & 5.364 & 4.844 & 3.611 & 1.34\\
200 & 15.809 & 14.546 & 12.912 & 1.13\\
\bottomrule\end{tabular}\end{center}
Complete sampler runs were warmed before timing; the following times are seconds.
\begin{center}\begin{tabular}{lrrrrr}\toprule
Mode & Seed & Original & Cached & Compiled & Speedup\\\midrule
pooled & 42 & 7.623 & 7.457 & 3.775 & 1.98\\
pooled & 123 & 7.849 & 7.756 & 3.936 & 1.97\\
hier & 42 & 11.806 & 11.566 & 7.684 & 1.51\\
hier & 123 & 11.909 & 11.602 & 7.492 & 1.55\\
\bottomrule\end{tabular}\end{center}
Two seeds and sequential process timings do not establish a precise universal
speedup. Hardware load, startup, batch size and model combination affect results.
First use compiles the kernel; cold first-evaluation timings are recorded
separately in the JSON metadata. The steady-state tables exclude compilation.
\section*{Usage and limitations}
With the project's optional Numba dependency installed, set
\texttt{settings.history\_backend: numba} in a copy of the copper YAML
configuration. Keep other settings unchanged. The main multi-calibration model
forwards this setting; separate Z blocks are unchanged. The constructor also
accepts the following option:
\begin{verbatim}
ModelMaterialStrength(..., history_backend="numba")
\end{verbatim}

The compiled driver requires finite float64 histories/state arrays, scalar
constants, and parameters broadcastable to one value per history. It checks
parameter validity and rejects nonfinite outputs. It does not promise identical
NumPy warnings, underflow exceptions, or partially updated state on invalid
inputs. Use the Python driver for custom models or unsupported inputs.
\section*{Tracking and reproduction}
Measured implementation: \texttt{ac87663}. Original commit: \texttt{6ccd5a1}.\\
Branch: \texttt{perf/compile-material-history}.\\
Step-1 checkpoint: \texttt{perf-step1-20260929}.\\
Step-2 results checkpoint: \texttt{perf-step2-20260929}.\\
All are local; nothing is published upstream.

Versioned documentation includes \texttt{docs/performance/STEP2.md}, the change
log, test log, compact results, profiles, artifact hashes, and this LaTeX source.
Full traces and raw timings are preserved in workspace
\texttt{benchmark/changes/step2/}. Previous artifacts and baseline fixtures
remain unchanged. The original checkout is in \texttt{impala-original/}.
\begin{verbatim}
.venv/bin/python benchmarks/compiled_material_history.py \
  --reference-root ../impala-original \
  --output ../benchmark/changes/step2-rerun
\end{verbatim}
Use a new output directory for every run. The benchmark fails on any exact
trajectory mismatch; it does not silently relax accuracy gates. Future changes
that alter arithmetic require numerical tolerances and multi-seed posterior,
predictive, R-hat and ESS-per-second comparisons before acceptance.
\end{document}
